\documentclass[11pt]{article}
\usepackage{preamble}
\setlength{\parskip}{4pt}

\begin{document}

\daybanner{AP Statistics \textperiodcentered\ Unit 5 \textperiodcentered\ Topic 5.4 \textperiodcentered\ B: 2027-02-11 \textperiodcentered\ E: 2027-04-07 \textperiodcentered\ TEACHER KEY}{One Line, Two Model Diagnostics}

\sectionbanner{CED TETHER}

{\small
\begin{itemize}[leftmargin=1.5em,itemsep=2pt,topsep=2pt]
  \item ENDURING UNDERSTANDING: DAT-1 Regression models may allow us to predict responses to changes in an explanatory variable.
  \item SKILLS:
  \item \textbf{Skill 2.B} --- Construct numerical or graphical representations of distributions.
  \item \textbf{Skill 2.A} --- Describe data presented numerically or graphically.
  \item LO DAT-1.E Represent differences between measured and predicted responses using residual plots. [Skill 2.B]
  \item EK DAT-1.E.1 The residual is the difference between the actual value and the predicted value: residual = y - $\widehat{y}$.
  \item EK DAT-1.E.2 A residual plot is a plot of residuals versus explanatory variable values or predicted response values.
  \item LO DAT-1.F Describe the form of association of bivariate data using residual plots. [Skill 2.A]
  \item EK DAT-1.F.1 Apparent randomness in a residual plot for a linear model is evidence of a linear form to the association between the variables.
  \item EK DAT-1.F.2 Residual plots can be used to investigate the appropriateness of a selected model.
\end{itemize}
}
\textbf{Essential question:} How should correlation and residuals be reconciled when judging a model for safety planning?\par
\textbf{DOK-3 skill:} 4.B \textperiodcentered\ \textbf{FRQ pattern:} {\small\texttt{correlation-\allowbreak{}residuals-\allowbreak{}model-\allowbreak{}choice-\allowbreak{}and-\allowbreak{}scope}}\par
\textbf{DOK rationale:} Part (c) weighs competing claims using the day's numerical or graphical evidence and explains a limitation or consequence; the judgment requires linked reasoning beyond a calculation.\par

\frameworkphaseheader{Do Now (first take) $\rightarrow$ Explore (video + follow-along) $\rightarrow$ Exit (finish + turn in)}{1 $\rightarrow$ 3}{45}{%
\begin{itemize}[leftmargin=1.5em,itemsep=2pt,topsep=2pt]
  \item Show the board at the bell and collect a one-sentence first take without confirming a claim.
  \item Run the named video follow-along, then allow about 10 minutes for parts (a)--(c).
  \item Ask for evidence linked to a claim. Collect the sheet and hand-score part (c) with E/P/I.
\end{itemize}
}{%
\begin{itemize}[leftmargin=1.5em,itemsep=2pt,topsep=2pt]
  \item Commit to a judgment, complete the follow-along, finish (a)--(c), revise, and turn in.
\end{itemize}
}{%
\begin{itemize}[leftmargin=1.5em,itemsep=2pt,topsep=2pt]
  \item Which number or visible feature supports your claim?
  \item Why does the other claim go beyond that evidence?
  \item What would you tell someone who chose the other claim?
\end{itemize}
}{Circulate and ask students to connect their choice to evidence. Score part (c) using the three required reasoning elements; do not require dropped extension work.}

\begin{callout}{\IconWarn\ WATCH FOR}{calloutred}
\begin{itemize}[leftmargin=1.5em,itemsep=2pt,topsep=2pt]
  \item A choice with copied numbers but no explanation of how they support it.
  \item Treating a model or average as a guaranteed individual outcome.
\end{itemize}
\end{callout}

\sectionbanner{SCORING PART (c) --- E / P / I}

\textbf{Expected elements}
\begin{itemize}[leftmargin=1.5em,itemsep=2pt,topsep=2pt]
  \item \textbf{reasoning-1} (required) --- Supports June by reconciling the large correlation with the U-shaped residuals
  \item \textbf{reasoning-2} (required) --- Links the 4-foot underprediction to a stopping-distance consequence
  \item \textbf{reasoning-3} (required) --- Requires a replacement residual check in the stated vehicle/surface setting
\end{itemize}

{\small\renewcommand{\arraystretch}{1.25}
\noindent\begin{tabular}{|p{0.5in}|p{5.9in}|}
\hline
\textbf{E} & Includes all three required reasoning elements above, with correct contextual evidence. \\ \hline
\textbf{P} & Includes at least one substantive correct reasoning link above, but omits or misstates another required element. \\ \hline
\textbf{I} & Includes no substantive correct reasoning link above; an unsupported choice or copied number alone is insufficient. \\ \hline
\end{tabular}}

\clearpage
\focusbanner{TODAY'S PROBLEM --- annotated}

Suppose a road-safety lab records braking distance $y$ at speeds $x$ from 10 through 50 mph on one vehicle and dry surface. The line is $\hat y=-10+2.5x$, the correlation is $r=0.997$, and the residual plot is shown. A residual is actual distance minus the line's prediction. At 50 mph, the actual distance is 119 feet.\par\smallskip \textbf{Ezra:} The large correlation and residuals on both sides of 0 justify using the line across the tested speeds.\par \textbf{June:} I would use where the residuals occur and the 50-mph error to judge safety, then check any replacement model separately.\par
\begin{center}
\begin{tikzpicture}
\begin{axis}[
  width=0.52\linewidth,
  height=1.40in,
  xmin=5, xmax=55, ymin=-5, ymax=5,
  xlabel={Speed (mph)}, ylabel={Residual (ft)}, grid=major,
  tick label style={font=\small}, label style={font=\small}
]
\addplot[only marks, mark=*, mark size=1.3pt, color=dayblue] coordinates {
  (10,4)
  (15,1)
  (20,-1)
  (25,-2)
  (30,-4)
  (35,-2)
  (40,-1)
  (45,1)
  (50,4)
};
\end{axis}
\end{tikzpicture}
\end{center}
\begin{firsttakebox}{FIRST TAKE (5 min)}
Would you trust this line to plan a safe stopping distance? Give one number or visible pattern in one sentence.\par
\answer{Not graded. Students should decide what each diagnostic establishes before judging the line.}
\end{firsttakebox}

\textit{Rules callout ``READING A RESIDUAL PLOT (keep on the page)'' is printed on the student sheet.}\par\medskip

\dokbadge{1}~\textbf{(a)} Read the residual at 50 mph. Does the line predict too high or too low there?\par
\answer{The residual is 4 feet, so the line predicts too low.}

\dokbadge{2}~\textbf{(b)} Verify the 50-mph residual using the equation. Describe where the residual plot is above and below zero.\par
\answer{At 50 mph, $\hat y=-10+2.5(50)=115$ feet, so the residual is $119-115=4$ feet. Residuals are positive at 10--15 and 45--50 mph, and negative at 20--40 mph, forming a U-shaped pattern.}

\dokbadge{3}~\textbf{(c)} Whose judgment is supported? Explain why the residual pattern matters even with the large correlation, and connect the 50-mph error to safety. State what you would check in a replacement model for this vehicle and dry surface. Write 3--4 sentences.\par
\answer{June is supported: $r=0.997$ shows a strong positive linear association, but the U-shaped residuals show a pattern of errors the line misses. Having both signs is not enough when the signs occur in a predictable order. At 50 mph the 4-foot underprediction could make a stopping zone too short. For this vehicle and dry surface, check that a replacement model's residuals have no clear pattern around zero before trusting it.}

\begin{summaryexitbox}{EXIT}
One thing the video changed about my first take: \hrulefill
\end{summaryexitbox}

\end{document}
